\begin{explanationbox}
\begin{minipage}{\linewidth}
\begin{center}\begin{center}
\begin{tikzpicture}[x=1cm,y=1cm,line cap=round,line join=round,font=\small]
\draw[black!25,rounded corners=3pt] (0,0) rectangle (14.4,9.2);
\node[font=\small\bfseries] at (7.2,8.7) {焦半径就是到焦点的距离：右支同侧短、异侧长};
\begin{scope}[shift={(4.5,4.3)},x=0.8cm,y=0.8cm]
\draw[black!30,->] (-4.4,0)--(4.5,0) node[right] {$x$};
\draw[black!30,->] (0,-3.8)--(0,4.1) node[above] {$y$};
\draw[black!65,thick,domain=-3.5:3.5,samples=100,variable=\z] plot ({-1*2*sqrt(1+\z*\z/5)},{\z});
\draw[black!65,thick,domain=-3.5:3.5,samples=100,variable=\z] plot ({1*2*sqrt(1+\z*\z/5)},{\z});
\coordinate (P) at (3,2.5);
\coordinate (Fone) at (-3,0);
\coordinate (Ftwo) at (3,0);
\draw[blue!70!black,line width=1.1pt] (Fone)--(P);
\draw[orange!85!black,line width=1.1pt] (Ftwo)--(P);
\foreach \q in {P,Fone,Ftwo} \fill (\q) circle (1.6pt);
\node[above left] at (P) {$P(3,\frac52)$};
\node[below left] at (Fone) {$F_1(-3,0)$};
\node[below right] at (Ftwo) {$F_2(3,0)$};
\path (Fone)--(P) node[pos=.34,above,sloped,blue!70!black] {$d_1=\frac{13}{2}$};
\node[right,orange!85!black] at (3,1) {$d_2=\frac52$};
\end{scope}
\node[align=left,anchor=north west] at (9.1,7.7) {$\dfrac{x^2}{4}-\dfrac{y^2}{5}=1$\\[1em]$a=2,\ c=3,\ e=\dfrac32$\\[1em]\textcolor{blue!70!black}{$d_1=ex_0+a$}\\[.7em]\textcolor{orange!85!black}{$d_2=ex_0-a$}\\[1em]$d_1-d_2=4=2a$};
\node at (7.2,.65) {$F_1$ 始终是左焦点，$F_2$ 始终是右焦点；焦点名称不随支改变。};
\end{tikzpicture}
\end{center}
\end{center}
先确认 $P$ 在双曲线上，再用横坐标正负区分两支。右支上两条焦半径均为正长度；左支须正确处理绝对值。
\end{minipage}
\par\medskip

\textbf{\textcolor{CExplanation}{一句话直觉}}

\textbf{两段距离相差固定的 $2a$；它们的中间值却是 $e|x_0|$。}
所以把两条焦半径直接看作“中间值减 $a$”与“中间值加 $a$”，
既能理解公式，又不用死记正负号。

\medskip
\textbf{\textcolor{CExplanation}{从熟悉的线段开始：先看右顶点}}

右顶点 $A(a,0)$ 与两焦点都在 $x$ 轴上，故
\[
 |AF_2|=c-a,\qquad |AF_1|=c+a.
\]
因为 $ea=c$，把 $c$ 写作 $ex_0$（此时 $x_0=a$），正好得到
“$ex_0-a$、$ex_0+a$”。这先解释了\textbf{为什么出现 $a$ 的加减}；
一般位置为什么仍成立，由后面的证明保证。

\medskip
\textbf{\textcolor{CExplanation}{最直观的拆解：$t-a$ 和 $t+a$}}

想象两条长度一长一短，差恰好 $2a$：
\[
 \underbrace{(t+a)}_{\text{较长}}-
 \underbrace{(t-a)}_{\text{较短}}=2a.
\]
这种设法\textbf{只有一个未知量 $t$}。
再用“两个距离的平方相减，纵坐标自动消失”就可求得
$t=e|x_0|$。\textbf{不需要先猜焦半径公式。}

\medskip
\textbf{\textcolor{CExplanation}{为什么纵坐标消失？}}

把 $P$ 关于 $x$ 轴翻折成 $P'(x_0,-y_0)$，
它到两焦点的距离完全不变；另外，
\[
 \begin{aligned}
 |PF_1|^2-|PF_2|^2
 &=\bigl[(x_0+c)^2+y_0^2\bigr]
   -\bigl[(x_0-c)^2+y_0^2\bigr]\\
 &=4cx_0.
 \end{aligned}
\]
两个 $y_0^2$ \textbf{直接抵消}，只剩下横坐标。

\medskip
\textbf{\textcolor{CExplanation}{如何做到不背左右支的四个符号？}}

先判 $P$ 在哪一支，再抓住“\textbf{同侧短，异侧长}”：
\begin{itemize}[leftmargin=*,itemsep=0.35em,topsep=0.35em]
 \item 右支 $x_0\ge a$：离右焦点 $F_2$ 近；
 $|PF_2|=ex_0-a$，$|PF_1|=ex_0+a$。
 \item 左支 $x_0\le-a$：离左焦点 $F_1$ 近；
 $|PF_1|=-ex_0-a$，$|PF_2|=-ex_0+a$。
 \item 两支合并：同侧短 $=e|x_0|-a$，异侧长 $=e|x_0|+a$。
\end{itemize}

\medskip
\textbf{\textcolor{CExplanation}{与椭圆放在一起理解}}

两者都来自\textbf{焦点距离的和/差为常数}：
\begin{itemize}[leftmargin=*,itemsep=0.3em,topsep=0.3em]
 \item 椭圆：两距离的\textbf{和}是 $2a$；适合设 $a+t$、$a-t$。
 \item 双曲线：两距离的\textbf{差的绝对值}是 $2a$；适合设 $t+a$、$t-a$。
\end{itemize}
这两个“一变量设法”只是把定值条件提前用掉了，运算负担很小。

\medskip
\textbf{\textcolor{CExplanation}{考场切入点}}

题目给出焦半径、焦点三角形边长、焦半径比例或 $x_0$ 时，
优先想到“\textbf{固定差 $2a$} + \textbf{同侧短、异侧长}”。
若给 $x_0$，直接代公式；若给两焦半径，先\textbf{加减}得到 $a$ 与 $e|x_0|$，
常比列两条开方距离式更快。

\end{explanationbox}
